% Gráficos canónicos de D3 en su orientación original.
\ProvideDocumentCommand\figuraHorizontalAqui{mmm}{\par\figuraHorizontal{#1}{#2}{#3}\clearpage}
% Contenido de D3 trasladado desde sus Markdown del Informe 2.
El modelo conserva el viaje ocurrido y sus respaldos sin mezclarlo con el costo estimado, el documento fiscal o el permiso de compartir datos. S3 fija reglas y alcance; S4.1 define servicios y contratos; S6/S7 organizan su construcción y migración; S9 comprueba que una decisión pueda reconstruirse a partir de la evidencia original. El CLIENTE recibe modelos por dominio, diccionario de atributos, gobierno de maestros, plazos de conservación y criterios de conciliación.

\section{Modelo}\label{sec:5-modelo}\label{sec:05-modelo}

Los seis contextos canónicos tienen propietarios explícitos. Compartir una identidad o referenciar un viaje no autoriza a un servicio a escribir en la base de otro. Las claves globales facilitan correlación y deduplicación; las restricciones de negocio se implementan en el contexto responsable.

\FloatBarrier
\subsection{Decisiones del modelo}\label{sec:5-entidades}\label{sec:05-decisiones-del-modelo}

La \tabref{05-tabla-5-1} vincula cada decisión con una necesidad del Caso, sin atribuir al legado un esquema que no se ha levantado.

\begin{tabla}[ancho=completo, fuente=condensada]{Decisiones de modelado y efecto operacional}{05-tabla-5-1}{L{36mm} X X}
Decisión & Necesidad & Consecuencia verificable \\ \midrule
Separar orden, viaje y asignación & Distinguir servicio solicitado, recursos reservados y recorrido efectivo & Reasignar recursos no borra la historia ni cancela automáticamente el viaje \\
Distinguir tracto y semirremolque & Aptitud y combinaciones de carga & Vigencias y compatibilidad por activo, no por una entidad genérica de camión \\
Evidencia, evaluación y permisos separados & Jornada impugnable y control del dato de terceros & Reconstruir qué fuentes se usaron, con qué regla y quién autorizó una excepción \\
Maestros con propietario único & Cuatro planillas de $\approx$6.000 vigencias y sistemas existentes & Identidad, respaldo y discrepancias conciliados antes de habilitar \\
Separar negocio, series y análisis & Asignación $\leq$30 s y ráfagas de reconexión & Ensayar cargas simultáneas y consultas de BI aisladas del OLTP \\
Versionar costo y liquidación & Combustible tardío y liquidación mensual & v1 $\leq$24 h con faltantes; posteriores versiones sin sobrescritura \\
\end{tabla}

La \tabref{05-tabla-5-1} fija la estructura que se comprueba en S9. El cifrado y los índices son controles complementarios; ninguno acredita por sí solo cumplimiento legal, exactitud documental o latencia.

\FloatBarrier
\subsection{Dominios de información y modelo conceptual}\label{sec:5-dominios}\label{sec:05-dominios-de-informacion-y-modelo-conceptual}

La arquitectura de datos se estructura en seis contextos delimitados con límites explícitos. En la \figref{05-audit-figura-5-1-erd} se despliega el modelo conceptual y relacional de persistencia de datos.

\figuraHorizontalAqui{figuras/05-datos/audit-figura-5.1-erd.pdf}{Modelo lógico de las relaciones críticas del viaje}{05-audit-figura-5-1-erd}

La \figref{05-audit-figura-5-1-erd} separa el ciclo del viaje de las asignaciones: una orden genera cero o varios viajes y cada viaje conserva cero o varias asignaciones históricas. Cada asignación confirmada identifica conductor, tracto y, cuando la carga lo exige, semirremolque; no puede solaparse con otra activa para esos recursos. Documento, costo y evaluación mantienen su identificador y relación con el viaje o la asignación. En el gráfico, EVAL\_\allowbreak{}ASIGNACION abrevia EVALUACION\_\allowbreak{}ASIGNACION y DOCUMENTO\_\allowbreak{}TRANSP. abrevia DOCUMENTO\_\allowbreak{}TRANSPORTE.

Los contextos gobiernan estos agregados: Planificación conserva OrdenTransporte y la reserva de Asignación; Operación conserva Viaje y su ejecución; Flota conserva Tractocamion y Semirremolque; Personas conserva Conductor, EvidenciaJornada y EvaluacionAsignacion; Liquidación conserva CostoVersion y Liquidacion. Transportista y Cliente tienen identidad compartida con un propietario de atributo declarado en MDM.

El ciclo propuesto del viaje es \texttt{PROGRAMADO $\rightarrow$ EN\_\allowbreak{}CARGA $\rightarrow$ EN\_\allowbreak{}RUTA $\rightarrow$ EN\_\allowbreak{}DESCARGA $\rightarrow$ CERRADO}, con \texttt{CANCELADO} como excepción trazable. \texttt{BLOQUEADA} es resultado de una evaluación de asignación, no un sexto estado que termine el viaje por fatiga: puede buscarse otro recurso y volver a evaluar. El rechazo de carga en destino conserva evento, motivo y resolución; no elimina recorrido ni costos ya incurridos. Un cierre de frontera registra demora y causa; su efecto sobre cobros depende de la regla contractual, sin congelar sobreestadías automáticamente por RT-10.05.

\FloatBarrier
\subsection{Estrategia de Gestión de Datos Maestros (MDM)}\label{sec:5-mdm}\label{sec:05-estrategia-de-gestion-de-datos-maestros-mdm}

En cumplimiento de \transv{RT-05.09}{11}, se implementa un marco centralizado de Master Data Management (MDM) con principio de Registro Maestro Único (\emph{Golden Record}), visualizado en la \figref{05-audit-figura-5-2-mdm}.

\figuraHorizontalAqui{figuras/05-datos/audit-figura-5.2-mdm.pdf}{Gestión de Datos Maestros (MDM) y Capa Anticorrupción}{05-audit-figura-5-2-mdm}

La \figref{05-audit-figura-5-2-mdm} identifica propietarios y conciliación: Personas gobierna conductor/evidencia/habilitación; Flota, activos y taller; Planificación, órdenes y puntos de servicio; el contable conserva identidad comercial y de facturación según el contrato de integración. Las tarifas operacionales cambian de autoridad por función en M16/M21 según S3; no existe doble escritura maestra.

Cada registro mantiene ID global y clave natural validada. El RUT usa validación sintáctica de dígito verificador y contraste documental; la patente se normaliza y contrasta con antecedentes autorizados. No se promete una API del Registro Civil ni una consulta oficial automática sin acceso acreditado. Los duplicados o fuentes contradictorias quedan en cuarentena con resolución del propietario y trazabilidad del origen.

\FloatBarrier
\subsection{Diccionario de datos y restricciones}\label{sec:5-diccionario}\label{sec:05-diccionario-de-datos-y-restricciones}

El archivo independiente \texttt{AUDIT-Subdocumento5-Anexos.pdf} entrega por atributo nombre, tipo, dominio, obligatoriedad, propietario y sensibilidad, conforme a \transv{RT-05.01}{11} y Comunicado 10 · capítulo 5. Las \figref{05-audit-figura-5-1-erd} y \figref{05-5-7-evidencia} muestran relaciones legibles; el cuerpo analiza las restricciones que determinan la operación.

Una vigencia referencia exactamente un conductor, un tracto o un semirremolque mediante tres FK nullable y un \texttt{CHECK} de un único sujeto. Se evita una FK polimórfica que PostgreSQL no pueda comprobar. El estado \texttt{VERIFICADO} requiere respaldo y fecha de verificación; también se comprueba vencimiento, tipo requerido y eventual revocación. Una fecha futura en una planilla no equivale a aptitud acreditada.

El viaje puede programarse sin recursos ni peso definitivo. La asignación sólo se confirma con recursos y evaluación válida; pasar a \texttt{EN\_\allowbreak{}RUTA} requiere además carga definitiva y documento conforme. La posición se almacena con latitud y longitud ambas presentes o ambas ausentes. Identidad, velocidad, odómetro y hora del evento se dejan nulos cuando la fuente no los proporciona; la hora de recepción y la procedencia siempre quedan registradas. No se generan ceros para representar desconocimiento.

\FloatBarrier
\subsection{Evidencia de jornada, cascada y cadena de custodia}\label{sec:5-evidencia-jornada}\label{sec:05-evidencia-de-jornada-cascada-y-cadena-de-custodia}

La evidencia se clasifica conforme al Subdocumento 3: fuente instrumental, nivel, estado, modalidad y veredicto son dimensiones distintas. La \tabref{05-tabla-5-2} conserva la correspondencia canónica.

\Needspace{20\baselineskip}
\begin{tabla}[ancho=completo, fuente=condensada]{Niveles de evidencia de jornada y uso en la validación}{05-tabla-5-2}{L{29mm} X X X}
\textbf{Origen o código conservado} & \textbf{Nivel canónico} & \textbf{Condición de interpretación} & \textbf{Modalidad y veredicto esperado} \\ \midrule
Tacógrafo digital (\texttt{TACO\_\allowbreak{}DIG}) & 1 & Original, identidad, vínculo e integridad & Evaluación legal completa \\
Equipo a bordo identificado (CAN/GPS) & 2 & Conductor y evento verificados & Modalidad completa \\
Telemetría de fábrica (\texttt{TEL\_\allowbreak{}FAB}) & 3 & Lectura autorizada y origen trazado & Sin inferir descanso personal \\
Reposo del camión (\texttt{REPOSO\_\allowbreak{}CAM}) & 4 & Estado del vehículo; sin coordenadas & Requiere nivel 5 en modalidad de datos \\
Atestación (\texttt{ATEST\_\allowbreak{}TRAN}) & 5 & Firma, responsable, fecha y viaje & Asignación con marca \\
Registro voluntario (\texttt{REG\_\allowbreak{}COND}) & 0 & Captura identificada & Informativo \\
Evento de portería & — & Entrada o salida & No habilitante \\
Sin evidencia (\texttt{SIN\_\allowbreak{}EVIDENCIA}) & \texttt{NULL} & Fuente ausente & Bloqueo; RN-04 acotada \\
Código desconocido & \texttt{NULL} & Original en cuarentena & Bloqueo hasta homologación \\
Evidencia inválida & \texttt{NULL} & Artefacto y motivo conservados & Bloqueo \\
\end{tabla}

La \tabref{05-tabla-5-2} resume categorías; las reglas completas no caben en una celda. El nivel clasifica la evidencia y no certifica por sí solo su valor judicial. Nivel 1 exige original, identidad y vehículo/viaje; nivel 2 requiere conductor identificado, no sólo el código instrumental. Nivel 3 conserva cobertura y hora de recepción, sin inferir descanso de una posición. Nivel 4 necesita atestación firmada de nivel 5 para asignar con marca en modalidad de datos; sin ella bloquea. El nivel 5 registra responsabilidad y no levanta bloqueos legales. Nivel 0 nunca habilita ni cambia el veredicto. Origen desconocido o inválido mantiene código/artefacto y motivo en cuarentena, sin renumerarse como nivel 0. RN-04 sólo cubre su excepción aprobada y no un incumplimiento legal. Los códigos de origen (CAN, GPS, portería u otro) se conservan sin convertirlos en niveles: CAN instalado en un equipo a bordo no equivale a telemetría de fábrica. Las plataformas de terceros se integran únicamente mediante lectura o exportación autorizada; no se presume API ni webhook.

La evaluación de asignación vinculada conserva por separado \texttt{modalidad\_\allowbreak{}adhesion}, \texttt{veredicto\_\allowbreak{}asignacion}, versión de regla y marca temporal, junto con las referencias a conductor, viaje y evidencia. La modalidad completa aplica la evidencia de nivel 2 prevista por S3; la modalidad de datos requiere reposo de nivel 4 complementado por atestación firmada de nivel 5 y produce asignación con marca; sin adhesión, la validación documental solo habilita con atestación válida y también deja marca. La falta de fuente bloquea, salvo la excepción acotada de RN-04. Un incumplimiento legal bloquea en todas las modalidades y el nivel 0 nunca modifica el resultado. Estos casos se cotejan con T-17, CP-SYS-05 y CP-UAT-09.

El origen recibido se conserva sin renumeración silenciosa. Una descarga de tacógrafo, un evento de equipo o una atestación mantiene original, emisor/dispositivo, versión de esquema, ID, marca temporal si existe, hora de recepción, hash y resultado de validación. \texttt{id\_\allowbreak{}conductor} puede ser nulo en un artefacto no identificado; ese artefacto se conserva para auditoría y no se convierte por sí solo en jornada personal habilitante. La vinculación exige identidad acreditada y período coherente.

\texttt{EVIDENCIA\_\allowbreak{}JORNADA} conserva artefactos válidos, desconocidos o inválidos; la ausencia se expresa en \texttt{EVALUACION\_\allowbreak{}ASIGNACION} como \texttt{SIN\_\allowbreak{}EVIDENCIA}, sin inventar una fila de evidencia. \texttt{nivel\_\allowbreak{}cascada} es NULL para origen desconocido/inválido; 0 sólo para \texttt{REG\_\allowbreak{}COND}; 1–5 según la fuente válida y la condición canónica. Un nivel válido no elimina un incumplimiento legal. La hora GNSS/NTP se compara con la de recepción y se registra deriva o calidad del reloj; no se llama reloj atómico ni se reemplaza la hora original para aparentar exactitud.

La \figref{05-5-2-custodia} recorre recepción, validación, sellado, conservación y rectificación.

\figuraHorizontalAqui{figuras/05-datos/5-2-custodia.pdf}{Cadena de custodia de evidencias y tratamiento de fuentes inválidas}{05-5-2-custodia}

El hash permite detectar alteración; no es firma electrónica ni vuelve inmutable una fila por sí solo. Se conserva el original firmado cuando procede, se encadenan registros con identificador de cadena y secuencia, y se anclan segmentos y manifiestos en almacenamiento WORM con retención bloqueada, bajo administración separada. Las cuentas de aplicación carecen de UPDATE/DELETE/TRUNCATE sobre evidencia. La protección frente al administrador se acredita en el almacenamiento inmutable y con verificación independiente; revocar permisos SQL al usuario aplicativo no restringe por sí solo a un superusuario.

Las rectificaciones generan un registro nuevo que referencia el original, sin sobrescribirlo. El primer hash previo es NULL y existe un único inicio por cadena; el orden es secuencia de custodia, distinto del orden de ocurrencia de eventos tardíos. Un acuse de sincronización requiere commit duradero. La confirmación probatoria incorpora el sellado y respaldo exigidos por el dominio; su falta se muestra y se ensaya, sin declarar evidencia plenamente custodiada cuando sólo fue admitida por el bus.

\begingroup\emergencystretch=2em
La \figref{05-5-7-evidencia} separa evaluación y evidencia mediante \textsf{EVALUACION\_\allowbreak{}EVIDENCIA}: una evaluación puede usar reposo de nivel 4 y atestación de nivel 5 simultáneamente. El veredicto siempre se registra, aun sin evidencia válida. Una excepción RN-04 conserva dos aprobadores, motivo, un viaje y revisión mensual; no levanta una causa legal. Las restricciones y nulabilidad se detallan en el anexo.\par\endgroup

\figuraHorizontalAqui{figuras/05-datos/5-7-evidencia.pdf}{Modelo lógico de evidencia, evaluación, excepción y permiso}{05-5-7-evidencia}

La \tabref{05-tabla-5-2} y la \figref{05-5-7-evidencia} se cotejan con CP-SYS-05/CP-UAT-09: nivel 0, nivel 5, modalidad de datos, atestación ausente, fuente desconocida y bloqueo legal. Una base migrada conserva los códigos anteriores, su tabla de correspondencia y la versión del mapeo; no altera la evidencia original para hacerla coincidir con el nuevo esquema.

\FloatBarrier
\subsection{Soberanía, permisos y bases de tratamiento}\label{sec:5-soberania}\label{sec:05-soberania-permisos-y-bases-de-tratamiento}

El permiso del transportista controla qué datos de sus recursos comparte, con quién, para qué finalidad y durante qué servicio (\caso{RT-16.30}{33}). No sustituye la base de tratamiento de datos personales del conductor. El modelo distingue permiso comercial, adhesión contractual y base jurídica; la localización es dato personal cuando identifica o permite identificar a una persona, sin clasificar automáticamente toda localización como categoría legal de dato sensible.

La \figref{05-5-3-permisos} muestra concesión, evaluación por petición, revocación y decisión de conservación. \texttt{PERMISO\_\allowbreak{}COMPARTIR} mantiene recurso, destinatario, alcance, finalidad, vigencia, versión y respaldo; \texttt{ACCESO\_\allowbreak{}DATO} conserva qué se mostró y bajo qué versión. Una lista JSON de clientes no reemplaza relaciones verificables ni la auditoría de cada acceso.

\figuraHorizontalAqui{figuras/05-datos/5-3-permisos.pdf}{Compartición autorizada, revocación y conservación}{05-5-3-permisos}

La revocación cesa nuevas consultas y emisiones cubiertas por el permiso, invalida caché/credenciales de acceso y propaga una versión al dispositivo. Durante falta de enlace no puede recibir una revocación nueva: el permiso local tiene vencimiento y sólo se transmiten datos tras comprobar la versión al reconectar. La transmisión de emergencia y los registros obligatorios siguen su base específica y minimización; revocar compartición comercial no borra evidencia de jornada ni elimina obligaciones de seguridad.

\art{4.2}{5} exige Ley 21.719 y, mientras esté vigente, Ley 19.628. La Ley 21.719 entra en vigor el 01-12-2026; el inicio supuesto de febrero de 2027 cae bajo el régimen nuevo. Curimón debe documentar finalidad, base, derechos, encargos, plazos y transferencias. El cifrado de campo exigido por RT-11.10 no reemplaza esos deberes (BCN, Ley 21.719 \parencite{d3_ley21719}).

Chile Central es primario y Brazil South es secundario propuesto en S4.3. Una réplica de datos personales es transferencia internacional: requiere identificar responsable y destinatario y documentar una vía aplicable del artículo 27 y las salvaguardias del artículo 28. Ubicación sudamericana y cifrado no acreditan por sí solos licitud. Si no se documenta el mecanismo, no se activa esa réplica; tampoco se declara satisfecho el DR hasta disponer de un sitio secundario admisible que supere RPO $\leq$15 min y RTO $\leq$4 h. San Bernardo puede sostener continuidad local ante una pérdida de enlace, pero no se acredita como recuperación frente a un desastre de la región primaria: comparte el área metropolitana y sus riesgos, conforme a RT-07.02. Mientras no exista una vía habilitante documentada para Brazil South o un sitio secundario admisible con independencia de fallos, la recuperación de esos datos queda sin acreditar. La revisión jurídica y la prueba de recuperación son condiciones de puesta en operación.

Una supresión solicitada se resuelve por dato y finalidad: conservar protegido lo sujeto a retención o litigio, cesar usos no autorizados y eliminar cuando proceda. El procedimiento incluye BD, originales, índices, cachés, exportaciones autorizadas, réplicas y respaldos; registra alcance, retenciones vigentes y fecha efectiva. NIST SP 800-88 Rev. 2 sustituye la revisión 1 desde septiembre de 2025. La destrucción criptográfica sólo se usa cuando el diseño de claves y todas sus copias lo permite; no se promete irrecuperabilidad absoluta ni purga instantánea de Key Vault (NIST, SP 800-88 Rev. 2 \parencite{d3_nist88r2}).

\section{Gestión de datos}\label{sec:5-gestion-datos}\label{sec:05-gestion-de-datos}

La gestión de datos de audIT define la gobernanza operacional del dato, los motores de persistencia justificados frente a los compromisos de disponibilidad y latencia, y las políticas de retención, respaldo y seguridad.

\FloatBarrier
\subsection{Motores, transacciones y decisión ante partición}\label{sec:5-motores}\label{sec:05-motores-transacciones-y-decision-ante-particion}

\transv{RT-05.02}{11} exige justificar por dominio consistencia y disponibilidad. CAP se aplica al comportamiento de un sistema distribuido durante una partición; ACID describe transacciones y no clasifica automáticamente un producto como CP/AP. La \tabref{05-tabla-5-3} y la \figref{05-audit-figura-5-3-cap} concretan la decisión operacional.

\begin{tabla}[ancho=completo, fuente=condensada]{Persistencia y conducta ante pérdida de enlace}{05-tabla-5-3}{L{36mm} X X}
Dominio & Motor y transacción & Decisión ante partición \\ \midrule
Asignación y vigencias & PostgreSQL 16 Flexible & Confirmación sólo con estado autoritativo \\
Captura a bordo & SQLite 3, WAL y outbox & Captura local y sincronización posterior \\
Telemetría & TimescaleDB separado; Event Hubs Premium & Búfer local y deduplicación al reconectar \\
Caché y sesiones & Azure Managed Redis & Revalidación o consulta autoritativa \\
Documentos y auditoría & Blob Storage inmutable y manifiestos & Custodia pendiente hasta acuse duradero \\
Analítica & ADLS Gen2, Delta y modelo dimensional & Actualización diferida con antigüedad publicada \\
\end{tabla}

La \tabref{05-tabla-5-3} mantiene una base de negocio consistente y una captura local que resiste particiones. Sin estado autoritativo no se confirma nueva reserva; RN-04 sólo cubre la indisponibilidad acotada de una fuente de evidencia. Una caída de nube permite capturar, no usar todas las funciones globales. TimescaleDB conserva semántica relacional y no se etiqueta como AP sólo por recibir eventos. Redis es dato derivado con versión y vencimiento: una copia obsoleta no habilita. El original local permanece cifrado hasta confirmar custodia, y un hash aislado no acredita inmutabilidad. En analítica se publican antigüedad, faltantes y versión. El documento fiscal permanece en el contable y requiere entrega conforme antes del movimiento; el búfer local no habilita emisión propia.

\figuraHorizontalAqui{figuras/05-datos/audit-figura-5.3-cap.pdf}{Consistencia del despacho y continuidad de captura durante una partición}{05-audit-figura-5-3-cap}

El protocolo Kafka de Event Hubs no implica un segundo clúster Apache Kafka ni transacciones exactamente una vez. Se usan outbox, entrega al menos una vez, orden por agregado y deduplicación duradera, conforme a S4.1.

\FloatBarrier
\subsection{Segregación arquitectónica OLTP versus analítica (Lakehouse)}\label{sec:5-oltp-olap}\label{sec:05-segregacion-arquitectonica-oltp-versus-analitica-lakehouse}

Para dar cumplimiento a \transv{RT-05.05}{11}, se separa la base transaccional de producción respecto del consumo analítico de reportería y cálculo de costo por kilómetro. La propagación usa captura de cambios (CDC Debezium) y Event Hubs Premium mediante protocolo Kafka. El objetivo propuesto de retraso CDC <60 s se verifica bajo carga antes de aceptarse; el flujo avanza hacia el repositorio analítico (Delta Lake), como esquematiza la \figref{05-audit-figura-5-4-cdc}.

\figuraHorizontalAqui{figuras/05-datos/audit-figura-5.4-cdc.pdf}{Segregación transaccional OLTP y analítica OLAP mediante CDC}{05-audit-figura-5-4-cdc}

El modelo de costeo por viaje (\caso{RT-05.29}{31}) administra el desfase inherente de las liquidaciones de combustible (hasta 40 días) y peajes mensuales, emitiendo una versión preliminar en no más de 24 horas y una versión consolidada definitiva post-cierre sin sobrescribir los registros originales.

\FloatBarrier
\subsection{Políticas de gobernanza, calidad y retención histórica}\label{sec:5-gobernanza}\label{sec:05-politicas-de-gobernanza-calidad-y-retencion-historica}

La calidad se mide en captura y migración según \transv{RT-05.04}{11} e ISO/IEC 25012. La \tabref{05-tabla-5-4} distingue los atributos que habilitan una transacción de los campos opcionales de una fuente.

\begin{tabla}[ancho=completo, fuente=condensada]{Calidad y método de comprobación}{05-tabla-5-4}{L{36mm} X X}
Dimensión & Objetivo de aceptación & Comprobación \\ \midrule
Completitud & 100 \% de campos exigidos para confirmar asignación/salida; reportar faltantes del resto & Numerador y denominador por entidad y fuente; no convertir NULL en cero \\
Exactitud & Identificadores válidos y documentos cotejados antes de habilitar & Validación sintáctica más respaldo autorizado; una API oficial no se presume \\
Consistencia & Sin asignaciones activas superpuestas ni vigencias sin sujeto/respaldo & FK, CHECK, unicidad/exclusión y ensayos concurrentes \\
Credibilidad & Origen, versión y horas de evento/recepción conservados; deriva visible & Contraste de reloj y orden de secuencia; cuarentena cuando no hay interpretación fiable \\
Accesibilidad & 99,9 \% mensual en servicios críticos durante operación & Indicadores de servicio extremo a extremo; HA del motor sola no acredita SLA \\
\end{tabla}

La \tabref{05-tabla-5-4} produce un tablero con errores por fuente, antigüedad, registros en cuarentena y resolución responsable. El 99,9 \% mensual se aplica a los servicios clasificados como críticos durante operación, igual que S4.2 y el Artículo 78.2; se calcula sobre minutos de transacción de negocio disponible / minutos exigibles del mes, por servicio, desde el punto de vista del usuario. No es disponibilidad aislada de PostgreSQL ni objetivo uniforme para servicios de otra criticidad. S9/T-17 deben medir el mismo perímetro y registrar exclusiones contractuales, si proceden. Las modificaciones y consultas sensibles conservan actor, dispositivo, valores anteriores/posteriores, motivo y correlación (RT-05.03 transversal y RT-16.09 del Caso).

La auditoría es append-only, encadenada y anclada en WORM bajo administración separada. La cuenta de negocio no modifica ni elimina registros de evidencia; una cuenta de administración de BD no puede modificar el original WORM. Los ensayos de integridad comprueban también omisión de eventos y eliminación de segmentos, con manifiestos firmados; el hash de una fila aislada no detecta por sí solo que esa fila desapareció.

La \tabref{05-tabla-5-5} recoge los mínimos de \caso{RT-05.10}{31}. No se atribuye a la ley tributaria una retención que aquí procede del contrato. El punto inicial del cómputo y las retenciones por litigio se documentan por dominio.

\begin{tabla}[ancho=completo, fuente=condensada]{Retención y recuperación por dominio}{05-tabla-5-5}{L{36mm} X X}
Dominio & Mínimo contractual & Preservación y objetivo \\ \midrule
Jornada y evidencia & 5 años & Original y rectificaciones; RTO $\leq$2 h \\
Documento y viaje & 6 años & Original fiscal y representación; RTO $\leq$2 h \\
Siniestros & 10 años & Expediente y litigio; RTO $\leq$4 h \\
Habilitaciones & Vigencia +5 años & Versiones y revocaciones; RTO $\leq$2 h \\
Carga peligrosa & 5 años & Documento y lista firmada; RTO $\leq$2 h \\
Llegada y salida & 3 años & Evento y origen; RTO $\leq$4 h \\
Liquidaciones & 6 años & Versiones y acuse contable; RTO $\leq$2 h \\
Posición y telemetría & 2 años en línea & Agregado posterior; RTO $\leq$4 h \\
\end{tabla}

La \tabref{05-tabla-5-5} conserva mínimos y diferencia objetivos internos de los contractuales. Jornada conserva evaluación, original y rectificaciones en segmentos sellados; documento vincula original fiscal y representación al viaje; siniestros retiene el expediente completo y mantiene litigio si procede. Las habilitaciones conservan respaldos, versiones y revocaciones. La telemetría pasa después de dos años a agregado diario por unidad/viaje para distancia y consumo, sin destruir evidencia sometida a otro plazo. Los RTO de dos horas son metas de diseño para los dominios indicados; el mínimo transversal de servicio crítico es RTO $\leq$4 h y RPO $\leq$15 min (\transv{RT-07.04}{17}). El desfase real de WAL, objetos y eventos debe medirse conjuntamente; un respaldo diario aislado no cumple RPO de quince minutos.

La política 3-2-1-1-0 incluye producción, copia recuperable localmente segregada y copia fuera de sitio, con almacenamiento de BD y objetos, al menos una copia inmutable y restauración sin errores verificada. Replicación continua no sustituye respaldo recuperable frente a corrupción. Se prueban restauración, conmutación y retorno al menos dos veces al año; el 0 es criterio de aceptación del ensayo, no resultado ya obtenido. La copia extranjera queda sujeta a 5.1.6 y S4.3. Al vencer retenciones y ausentes litigios, la sanitización se documenta con eliminación efectiva de todas las copias o expiración controlada de respaldos, sin reintroducir datos suprimidos al restaurar.

\section{Estrategia de migración}\label{sec:5-migracion}\label{sec:05-estrategia-de-migracion}

Transportes Curimón S.A. administra cerca de 6.000 vigencias dispersas en cuatro planillas sin validación referencial. En cumplimiento de \caso{RT-05.15}{31}, ningún registro migrado se considerará habilitante para despachar sin su correspondiente verificación documental individual. En la \tabref{05-tabla-5-6} se detalla el alcance cuantitativo y en la \figref{05-audit-figura-5-5-migracion} el flujo metodológico en cuatro fases.

\Needspace{20\baselineskip}
\begin{tabla}[ancho=completo, fuente=condensada]{Alcance de migración histórica del Caso 10}{05-tabla-5-6}{L{36mm} X X}
\textbf{Dominio Histórico a Migrar} & \textbf{Volumen del Caso 10} & \textbf{Criterio de Aceptación y Conciliación} \\ \midrule
Flota y semirremolques & 374 tractos; 210 semirremolques propios & Inventario y respaldo 1:1 \\
Conductores y transportistas & 454 choferes; 148 dueños & RUT y contrato \\
Clientes y geocercas & 84 clientes; 1.400 puntos & Dirección y polígono \\
Vigencias vivas & $\approx$6.000; cuatro planillas & Respaldo individual \\
Viajes históricos & Cinco años; $\approx$480.000 estimados & Totales, código y fecha \\
Liquidaciones históricas & Seis años; $\approx$10.656 estimadas & Importe, moneda y redondeo \\
Siniestros históricos & Total disponible; conteo pendiente & Expediente íntegro \\
\end{tabla}

La \tabref{05-tabla-5-6} distingue alcance obligatorio de estimaciones derivadas: los terceros de flota se migran según inventario real, no como 210 semirremolques propios. En maestros se exige conciliación 1:1 contra inventario y respaldo autorizado, RUT validado, contrato vigente y geocerca comprobada. Cada vigencia debe verificarse documentalmente; el hash SHA-256 comprueba integridad, no validez del documento. Viajes se concilian por totales, código y fecha; liquidaciones por importe, moneda y redondeo contra fuente y contable; siniestros por integridad de expediente y peritaje. El Caso aporta una tasa anual actual, no el número real de filas de los últimos cinco/seis años. En el levantamiento se cuentan registros, adjuntos y bytes, incluidos históricos sin persona/vehículo homologado; el archivo original se conserva en cuarentena hasta resolverlos.

La \figref{05-audit-figura-5-5-migracion} muestra cuatro fases: extracción/perfilamiento con copia y hash del origen; homologación y deduplicación con mapeos aprobados; acreditación de cada vigencia y custodia del respaldo; ensayos en PREPROD, conciliación y carga final de diferencias. La ventana se estima con conteos/bytes reales y tasas de los ensayos; se concilia con S7 antes del corte. Los paquetes de datos EDT 7.x y cortes por función M16/M21 gobiernan la ejecución.

La aceptación exige conciliación de totales, referencias y documentos, ausencia de vigencias habilitantes sin respaldo y cuadratura de liquidaciones con diferencias justificadas. Toda inconsistencia se identifica por registro, propietario y resolución, sin afirmar que una carga masiva elimina errores. Se prueban reversa, delta final e idempotencia; nunca se borran documentos fiscales ya emitidos ni evidencias válidas generadas durante la convivencia. Los cortes respetan los congelamientos y la implantación por terminal/unidad de S7 (\caso{RT-10.05}{32}).

\figuraHorizontalAqui{figuras/05-datos/audit-figura-5.5-migracion.pdf}{Metodología de extracción, saneamiento, carga y conciliación histórica}{05-audit-figura-5-5-migracion}

La \figref{05-audit-figura-5-5-migracion} muestra las puertas de acreditación y conciliación. Un registro en cuarentena no habilita despacho; el corte requiere acta y aceptación de los casos vinculados en T-17, con discrepancias resueltas o excluidas explícitamente del uso habilitante.

\section{Estrategia de desempeño}\label{sec:5-desempeno}\label{sec:05-estrategia-de-desempeno}

La meta contractual de asignación es $\leq$30 s extremo a extremo (\caso{RT-09.01}{32}). Menos de 2 s es objetivo interno de un caso nominal, no medición ni garantía derivada de usar caché. La \figref{05-audit-figura-5-6-desempeno} muestra que una proyección sólo habilita cuando su versión y actualidad se validan contra el estado autoritativo; si hay conflicto se reevalúa o rechaza dentro del plazo total.

\figuraHorizontalAqui{figuras/05-datos/audit-figura-5.6-desempeno.pdf}{Estado autoritativo, caché e índices especializados}{05-audit-figura-5-6-desempeno}

La \figref{05-audit-figura-5-6-desempeno} separa consultas puntuales del histórico y la analítica. Se proponen B-tree sobre IDs/FK y rangos de asignación, GiST en geocercas y BRIN sólo donde el orden físico temporal lo justifique. Un RUT cifrado no se indexa como texto en claro: la deduplicación usa un índice ciego HMAC con clave separada, alcance documentado y acceso restringido; se evalúa el riesgo de correlación. Las búsquedas parciales no se ofrecen sobre ciphertext como si fueran equivalentes a texto.

Con $\approx$96.000 viajes/año, \texttt{VIAJE} se mantiene sin partición inicialmente para conservar PK/FK simples. Se particionan series, eventos y auditoría por tiempo; si luego se particiona negocio, PostgreSQL exige incluir la clave de partición en la restricción única aplicable y propagarla a las referencias. No se propone una PK global por UUID incompatible con la partición mensual. La extensión TimescaleDB en Azure es Apache-2; particionado y agregaciones usan sólo funciones verificadas en esa edición, sin asumir compresión comercial.

Redis conserva versiones, vencimiento y permisos. Revocar una habilitación o autorización invalida la proyección; un TTL aislado no asegura actualidad. El caché no es el único registro de idempotencia. Vistas y consultas de preliquidación/BI residen en la capa analítica, con refresco y marca de antigüedad; no compiten con la base de despacho.

S4.1.20 deriva \textasciitilde{}119,82 millones de posiciones/año en el escenario completo y distingue las 61 fuentes iniciales de motor del resto. El anexo de S5 estima bytes de negocio, evidencia y migración con tamaños supuestos. La aceptación mide máximos y percentiles del despacho, retraso de CDC, consumo de WAL, CPU/IO, reconstrucción de evidencia y confirmación de cada camión. Se separan 380 sesiones nominales, 570 en estrés, crecimiento 3$\times$ y 300 reconexiones: no son la misma prueba ni equivalen a TPS. La línea base concurrente coincide con el perfil de carga del Subdocumento 4: 380 sesiones nominales y 570 en estrés.
